% Part: computability
% Chapter: computability-theory
% Section: non-comp-set

\documentclass[../../../include/open-logic-section]{subfiles}

\begin{document}

\olfileid{cmp}{thy}{ncp}
\olsection{గణనీయంకాని సమితులు ఉన్నాయి}


ప్రతి గణనీయ సమితీ గణనీయంగా లెక్కించదగినదని పైన చూశాం. దాని
విలోమం నిజమా? సాధారణంగా నిజం కాదని కింది సిద్ధాంతం చూపుతుంది.

\begin{thm}\ollabel{thm:K-0}
$K_0$ అనేది $\Setabs{\tuple{e, x}}{\cfind{e}(x) \fdefined}$ సమితి
అనుకుందాం. అప్పుడు $K_0$ గణనీయంగా లెక్కించదగినది, కాని గణనీయం కాదు.
\end{thm}

\begin{proof}
$K_0$ గణనీయంగా లెక్కించదగినదని చూడటానికి, కింది విధంగా నిర్వచించిన
ప్రమేయం~$f$ యొక్క నిర్వచన క్షేత్రమే అది అని గమనించండి:
\[
f(z) = \umin{y}{(\len{z} = 2 \land T((z)_0, (z)_1, y))}.
\]
ఎందుకంటే, $\cfind{e}(x)$ నిర్వచితమైతే
$f(\tuple{e, x})$ ఆగే గణనా అనుక్రమాన్ని కనుగొంటుంది;
$\cfind{e}(x)$ అనిర్వచితమైతే $f(\tuple{e, x})$ కూడా అనిర్వచితమే;
అంతేకాక $z$ ఒక జతను కూడా సంకేతీకరించకపోతే $f(z)$ కూడా అనిర్వచితం.

$K_0$ గణనీయం కాదనే వాస్తవం, ఆగే సమస్య అనిర్ణయనీయమనే
\olref[hlt]{thm:halting-problem} ఫలితమే.
\end{proof}

$K_0$ అనేది $\cfind{e}(x) \fdefined$ అయ్యే జతల
$\tuple{e,x}$ సమితి; అంటే $\cfind{e}$ ఇన్‌పుట్~$x$పై నిర్వచితమై
(ఆగి) ఉంటేనే $\tuple{e,x} \in K_0$. అందువల్ల దీన్ని ``ఆగే
సమితి'' అని కూడా అంటారు. $K = \Setabs{e}{\cfind{e}(e) \fdefined}$
అనేది ``స్వయంగా ఆగే సమితి.'' దీన్ని తరచుగా ఒక ప్రామాణిక
అనిర్ణయనీయ సమితిగా వాడతారు.

\begin{thm}\ollabel{thm:K}
స్వయంగా ఆగే సమితి $K = \Setabs{e}{\cfind{e}(e) \fdefined}$
\tetoken{గణనీయంగా లెక్కించదగినది}{!!{c.e.}}, కాని నిర్ణయించదగినది కాదు.
\end{thm}

\begin{proof}
  $K$ నిర్ణయించదగినదని, అంటే దాని లక్షణ ప్రమేయం $\Char{K}$
  గణనీయమని అనుకుందాం. ఇలా ఉంచండి:
  \[d(e) = \begin{cases}
  1 & \text{$\Char{K}(e) = 0$ అయితే}\\
  \fundefined & \text{ఇతర సందర్భాల్లో.}
  \end{cases}
  \]
  $k$ అనేది $d$ యొక్క సూచిక అనుకుందాం; అంటే $d \simeq
  \cfind{k}$. అప్పుడు $d(k) \simeq \cfind{k}(k)$. కానీ $d$
  నిర్వచనం నుంచి $d(k) \fdefined$ అయితేనే $\cfind{k}(k)
  \fundefined$ వస్తుంది; ఇది వైరుధ్యం.

  $K$ అనేది $f(x) = \umin{y}{T(x,x,y)}$ యొక్క నిర్వచన క్షేత్రం;
  కనుక అది \tetoken{గణనీయంగా లెక్కించదగినది}{!!{c.e.}}.
\end{proof}

\end{document}
